\thispagestyle{empty}
\hypertarget{frontmatter-page-006}{}
\null
\clearpage

\hypertarget{frontmatter-page-007}{}
\section*{前言}
\label{front:preface}
\addcontentsline{toc}{section}{前言}

Fourier 分析是数学的一个庞大分支, 其出发点是研究 Fourier 级数和 Fourier 积分. 然而, 它包含多种观点和技巧, 因而可以有许多题为 ``Fourier 分析'' 的不同导论. 本书的目标是研究 A. P. Calderón 与 A. Zygmund 在 20 世纪 50 年代引入 Fourier 分析的实变量方法.

我们在第 \ref{chap:fourier-series-integrals} 章先回顾 Fourier 级数和 Fourier 积分, 然后在第 \ref{chap:hardy-littlewood-maximal-function} 章和第 \ref{chap:hilbert-transform} 章介绍该领域的两个基本算子: Hardy--Littlewood 极大函数和 Hilbert 变换. 尽管它们早于 Calderón 与 Zygmund 的方法出现, 我们仍从他们的观点来处理这些算子. 这些方法旨在研究 Hilbert 变换在高维中的类似物; 它们在应用中十分重要. 这类算子称为奇异积分, 我们在第 \ref{chap:singular-integrals-one} 章和第 \ref{chap:singular-integrals-two} 章中讨论它们及其现代推广. 随后我们考察 20 世纪 70 年代出现的众多成果中的两项. 第 \ref{chap:h1-and-bmo} 章研究 $H^1$、BMO 与奇异积分之间的关系, 第 \ref{chap:weighted-inequalities} 章介绍加权范数不等式的初等理论. 第 \ref{chap:littlewood-paley-multipliers} 章讨论 Littlewood--Paley 理论; 它起源于 20 世纪 30 年代, 此后得到广泛发展, 并包含许多应用. 本章给出的应用可用于研究 Fourier 乘子, 其中也会使用加权不等式理论. 全书以 20 世纪 80 年代的一项重要结果结束, 即所谓的 $T1$ 定理; 它对这一领域具有关键意义.

每章末尾都有一节, 其中我们尽量介绍正文未讨论的进一步结果, 并为有兴趣的读者提供参考文献. 许多书籍以及所有引用的论文只出现在这些注记中; 书末参考文献则只收录深入讨论本书所述思想的著作.

\clearpage
\hypertarget{frontmatter-page-008}{}

本书材料来自 1988--89 学年在 Universidad Autónoma de Madrid 讲授的一门研究生课程. 其中一部分以我 1985 年秋在同一所大学作为学生听 José Luis Rubio de Francia 授课时所记的笔记为基础. 他似乎曾打算把课程整理成文, 但英年早逝使他未能完成. 因此, 我冒昧地使用了他的思想来写作本书; 这些思想既来自课堂, 也来自我们在走廊和黑板前许多愉快的交谈. 尽管本书献给他, 我几乎把它看作一部合作著作. 我还要感谢 Universidad Autónoma de Madrid 的朋友们, 是他们鼓励我讲授这门课程并写作本书.

本书最初以西班牙文收入 Universidad Autónoma de Madrid 的 \emph{Colección de Estudios} 出版, 1991; 随后由 Addison--Wesley/Universidad Autónoma de Madrid 于 1995 年联合再版, 只作了少量排印订正. 从一开始, 就有一些同事认为英文译本会引起兴趣, 但我从未着手翻译. David Cruz-Uribe 教授提出翻译本书时, 我立即接受. 我马上意识到, 文本不能保持原样, 因为过去十年的许多进展需要加入各章结尾的资料性部分, 同时还要加入初版遗漏的若干主题. 因此, 尽管本书核心只作了少量修改, 题为 ``注记与进一步结果'' 的各节却大幅扩充, 纳入了新的主题、结果和参考文献.

如果没有 Cruz-Uribe 教授极为宝贵的贡献, 本书的更新工作不可能取得现在的成果. 除了通读全文、建议修改和澄清晦涩之处外, 他还出色地扩充了上述注记, 查找参考文献, 并提出应当收入的新结果. 与原版相比, 本书的改进无疑是我们共同工作的成果. 我非常感谢他与我分享远远超出普通译者职责范围的学科知识.

\begin{flushright}
Javier Duoandikoetxea\\
Bilbao, June 2000
\end{flushright}

\clearpage
\hypertarget{frontmatter-page-009}{}

\noindent\textit{致谢:} 译者感谢 Ford Foundation 和 Trinity College 的 Dean of Faculty 在 1998--99 学年给予的慷慨支持. 正是在这一整年的学术休假期间, 本项目得以构思, 并完成了译稿初稿.

\clearpage
\thispagestyle{empty}
\hypertarget{frontmatter-page-010}{}
\null
\clearpage

\hypertarget{frontmatter-page-011}{}
\section*{预备知识}
\label{front:preliminaries}
\addcontentsline{toc}{section}{预备知识}

这里回顾一些记号和基本结果, 但我们假定读者大都已经熟悉. 更多资料例如见 Rudin \MBUnresolvedReference{citation}{[14]}.

通常我们在 $\mathbb R^n$ 中工作. Euclidean 范数记作 $|\cdot|$. 若 $x\in\mathbb R^n$ 且 $r>0$, 则
\[
B(x,r)=\{y\in\mathbb R^n:|x-y|<r\}
\]
是以 $x$ 为中心、$r$ 为半径的球. $\mathbb R^n$ 上的 Lebesgue 测度记作 $dx$, $\mathbb R^n$ 中单位球面 $S^{n-1}$ 上的测度记作 $d\sigma$. 若 $E\subset\mathbb R^n$, 则 $|E|$ 表示其 Lebesgue 测度, $\chi_E$ 表示其示性函数: 当 $x\in E$ 时 $\chi_E(x)=1$, 当 $x\notin E$ 时为 $0$. ``几乎处处'' 或 ``对几乎每个 $x$'' 是指相应性质在一个零测集以外成立; 分别简写为 ``a.e.'' 和 ``a.e. $x$''.

若 $\alpha=(\alpha_1,\ldots,\alpha_n)\in\mathbb N^n$ 是多重指标, 且 $f:\mathbb R^n\to\mathbb C$, 则
\[
D^\alpha f=
\frac{\partial^{|\alpha|}f}{\partial x_1^{\alpha_1}\cdots\partial x_n^{\alpha_n}},
\]
其中 $|\alpha|=\alpha_1+\cdots+\alpha_n$, 且 $x^\alpha=x_1^{\alpha_1}\cdots x_n^{\alpha_n}$.

设 $(X,\mu)$ 是测度空间. $L^p(X,\mu)$, $1\leq p<\infty$, 表示从 $X$ 到 $\mathbb C$ 的、$p$ 次幂可积的函数所成的 Banach 空间; $f\in L^p(X,\mu)$ 的范数为
\[
\|f\|_p=\left(\int_X|f|^p\,d\mu\right)^{1/p}.
\]

\clearpage
\hypertarget{frontmatter-page-012}{}

$L^\infty(X,\mu)$ 表示从 $X$ 到 $\mathbb C$ 的本质有界函数所成的 Banach 空间; 更确切地说, 对某个 $C>0$ 满足
\[
\mu(\{x\in X:|f(x)|>C\})=0
\]
的函数 $f$. $f$ 的范数 $\|f\|_\infty$ 是所有具有这一性质的常数的下确界. 通常 $X$ 是 $\mathbb R^n$ 或其子集, 且 $d\mu=dx$; 此时常常不写测度或空间, 而只写 $L^p$. 对一般测度空间, 我们也常用 $L^p(X)$ 代替 $L^p(X,\mu)$; 若 $\mu$ 绝对连续且 $d\mu=w\,dx$, 则写作 $L^p(w)$. $p$ 的共轭指数始终记作 $p'$:
\[
\frac1p+\frac1{p'}=1.
\]

$L^p$ 上的三角不等式有一个积分版本, 称为 Minkowski 积分不等式, 本书将反复使用它. 给定具有 $\sigma$-有限测度的测度空间 $(X,\mu)$ 和 $(Y,\nu)$, 该不等式为
\[
\left(\int_X\left|\int_Y f(x,y)\,d\nu(y)\right|^p d\mu(x)\right)^{1/p}
\leq\int_Y\left(\int_X|f(x,y)|^p\,d\mu(x)\right)^{1/p}d\nu(y).
\]

$\mathbb R^n$ 上两个函数 $f,g$ 的卷积定义为
\[
f*g(x)=\int_{\mathbb R^n}f(y)g(x-y)\,dy
=\int_{\mathbb R^n}f(x-y)g(y)\,dy,
\]
只要该表达式有意义.

测试函数空间是 $C_c^\infty(\mathbb R^n)$, 即紧支的无穷可微函数空间, 以及所谓的 Schwartz 函数空间 $\mathcal S(\mathbb R^n)$. Schwartz 函数是无穷可微且在无穷远处迅速衰减的函数; 更确切地说, 函数及其所有导数的衰减都快于任意多项式的增长. 赋予适当拓扑后, 它们的对偶分别是分布空间和缓增分布空间. 分布与测试函数的卷积可以定义如下: 若 $T\in C_c^\infty(\mathbb R^n)'$ 且 $f\in C_c^\infty(\mathbb R^n)$, 则
\[
T*f(x)=\langle T,\tau_x\widetilde f\rangle,
\]
其中 $\widetilde f(y)=f(-y)$, 且 $\tau_xf(y)=f(x+y)$. 若 $T$ 是局部可积函数, 该定义与前面的定义一致. 类似地, 可以取 $T\in\mathcal S(\mathbb R^n)'$ 和 $f\in\mathcal S(\mathbb R^n)$. 对偶配对不加区别地记作 $\langle T,f\rangle$ 或 $T(f)$.

方括号中的引用指向书末参考文献中的条目.

最后说明, $C$ 表示正的常数; 即使在同一串不等式中, 它也可能在不同位置取不同的值.
